% Part: many-valued-logic
% Chapter: syntax-and-semantics
% Section: valuations-sat

\documentclass[../../../include/open-logic-section]{subfiles}

\begin{document}

\olfileid[fr]{mvl}{syn}{val}

\olsection{Valuations et satisfaction}

\begin{defn}[Valuations]
Soit~$V$ un ensemble de valeurs de
vérité. Une \emph{valuation} du
langage~$\Lang{L}$ dans~$V$ est
une fonction~$\pAssign{v}$ qui
attribue à chaque variable
propositionnelle une valeur de~$V$ :
$\pAssign{v}\colon\PVar\to V$.
\end{defn}

\begin{defn}\ollabel{defn:pValue}
Étant donnée une valuation
$\pAssign{v}$ dans l'ensemble~$V$
des valeurs de vérité d'une logique
plurivalente~$\Log L$, définissons
par induction la fonction d'évaluation
$\pValue{v}\colon\Frm[L]\to V$ :
\begin{enumerate}
\item $\pValue{v}(\Obj p_n)
  =\pAssign{v}(\Obj p_n)$;
\item si $\star$ est un connecteur
  d'arité~$0$, alors
  $\pValue{v}(\star)
    =\tf{\star}[\Log L]$;
\item si $\star$ est un connecteur
  d'arité~$n\geq1$, alors
  \[
    \pValue{v}(\star(!A_1,\dots,!A_n))
      =\tf{\star}[\Log L]
        (\pValue{v}(!A_1),\dots,
         \pValue{v}(!A_n)).
  \]
\end{enumerate}
\end{defn}

\begin{defn}[Satisfaction]
\ollabel{defn:satisfaction}
La formule~$!A$ est \emph{satisfaite}
par la valuation~$\pAssign{v}$
dans la logique~$\Log L$,
$\pSat{v}{!A}[\Log L]$, si et
seulement si
$\pValue{v}(!A)[\Log L]\in V^+$,
où~$V^+$ est l'ensemble des valeurs
de vérité désignées de~$\Log L$.

La notation $\pSat/{v}{!A}[\Log L]$
signifie que $\pSat{v}{!A}[\Log L]$
ne vaut pas. Si~$\Gamma$ est un
ensemble de formules, alors
$\pSat{v}{\Gamma}[\Log L]$ si et
seulement si $\pSat{v}{!A}[\Log L]$
pour toute formule~$!A\in\Gamma$.
\end{defn}

\end{document}
